\scanpage{28}
\mspage{13}
où $s_0,s_1$ sont resp. les nbs de sommets et d'arêtes du graphe $G$. Or
\[
-3s_0+2s_1=-3(\underbrace{s_0-s_1}_{1-h_1})-s_1=-3+3h_1-s_1=3(h_1-1)-s_1
\]
d'où
\[
\dim\mathcal M_G/\mathbb Z=3(\underbrace{(\sum g_\alpha)+h_1}_{g}-1)+\nu-s_1
\]
et enfin la formule
\[
(\unsure{28})\qquad\dim\mathcal M_G/\mathbb Z=\underbrace{[3(g-1)+\nu]}_{=\dim\widehat{\mathcal M}_{g,\nu}/\mathbb Z}-s_1,
\]
i.e.
\[
(\unsure{29})\qquad\operatorname{codim}_{\widehat{\mathcal M}_{g,\nu}}\mathcal M_{[G]}=s_1
\]
(nb d'arêtes du MD-graphe $G$, i.e. ``nb de pts doubles'' des courbes \unsure{standard|standart} de type $G$).

L'entier (27) est appelé \uline{dimension modulaire} du MD-graphe $G$
\[
30)\qquad\begin{aligned}
\mu(G)&=\sum_{\alpha\in S}(3g_\alpha-1)+\widehat\nu_\alpha)=3\sum g_\alpha-3s_0+\nu+2s_1\\
&=(3(g-1)+\nu)-s_1,
\end{aligned}
\]
l'entier $s_1=\operatorname{card}A=(\operatorname{card}\widetilde A)/2$ s'appelle la codimension modulaire. Elle est nulle sss $G$ est un MD-graphe trivial (réduit à un pt isolé, sans arêtes, mais
\mspage{27}
